% VI33 detailed challenge solutions -- VI/24+ proof-depth standard.

% VI33-DETAILED-CHALLENGE-01
\subsection*{Challenge VI.33.C1 solution}
\begin{solution}
Let \(X\) be a scheme and let
\[
\mathcal I\subseteq\mathcal O_X
\]
be a quasi-coherent ideal sheaf.  We construct a closed subscheme of \(X\) affine-locally and then glue the local
pieces.

Choose an affine open cover
\[
X=\bigcup_iU_i,
\qquad
U_i=\Spec A_i.
\]
Because \(\mathcal I\) is quasi-coherent,
\[
\mathcal I|_{U_i}\cong\widetilde{I_i}
\]
for some \(A_i\)-module \(I_i\).  The inclusion
\[
\mathcal I|_{U_i}\hookrightarrow\mathcal O_{U_i}\cong\widetilde{A_i}
\]
corresponds, under the affine equivalence, to an injective \(A_i\)-linear map
\[
I_i\hookrightarrow A_i.
\]
Compatibility with multiplication by sections of \(\mathcal O_{U_i}\) shows that its image is an ideal, which we
again denote by \(I_i\subseteq A_i\).

Define
\[
Z_i=\Spec(A_i/I_i)
\hookrightarrow
U_i.
\]
On an overlap \(U_i\cap U_j\), refine by distinguished affine opens.  If
\[
D(f)\subseteq U_i
\]
lies in the overlap, then quasi-coherence gives
\[
\mathcal I(D(f))
\cong
(I_i)_f.
\]
The same restricted ideal sheaf computed from \(U_j\) gives the corresponding localization of \(I_j\). Therefore
the closed subschemes
\[
Z_i|_{U_i\cap U_j}
\quad\text{and}\quad
Z_j|_{U_i\cap U_j}
\]
are canonically isomorphic: both are cut out by the same restricted ideal sheaf
\(\mathcal I|_{U_i\cap U_j}\).

On triple overlaps these identifications satisfy the cocycle condition because they are induced by equality of the
same subsheaf of \(\mathcal O_X\).  Hence the \(Z_i\) glue to a scheme \(Z\), together with a morphism
\[
i:Z\hookrightarrow X
\]
which is affine-locally a closed immersion and therefore is a closed immersion.

Moreover the kernel of
\[
\mathcal O_X\longrightarrow i_*\mathcal O_Z
\]
is exactly \(\mathcal I\), since this is true on every affine \(U_i\):
\[
0\to\widetilde{I_i}
\to\mathcal O_{U_i}
\to(i_*\mathcal O_Z)|_{U_i}
\to0.
\]

Conversely, if
\[
i:Z\hookrightarrow X
\]
is a closed immersion, then on every affine open \(U=\Spec A\subseteq X\),
\[
Z\cap U=\Spec(A/I)
\]
for some ideal \(I\subseteq A\), and the kernel ideal sheaf restricts to \(\widetilde I\). Hence that ideal sheaf is
quasi-coherent.

Therefore
\[
\boxed{
\{\text{quasi-coherent ideal sheaves }\mathcal I\subseteq\mathcal O_X\}
\longleftrightarrow
\{\text{closed subschemes of }X\}.
}
\]
The correspondence sends \(\mathcal I\) to the quotient locally ringed space defined by
\(\mathcal O_X/\mathcal I\).
\end{solution}

% VI33-DETAILED-CHALLENGE-02
\subsection*{Challenge VI.33.C2 solution}
\begin{solution}
Let
\[
X=\Spec A.
\]
The affine equivalence already identifies all quasi-coherent sheaves with \(A\)-modules.  Restrict it to the
finiteness classes on the two sides.

Suppose first that \(P\) is a finitely generated projective \(A\)-module.  For each
\(\mathfrak p\in\Spec A\),
\[
P_{\mathfrak p}
\]
is a finite projective module over the local ring \(A_{\mathfrak p}\), hence free of some finite rank.  A basis
extends to a distinguished neighborhood \(D(f)\ni\mathfrak p\), so
\[
P_f\cong A_f^{\oplus r}.
\]
Thus
\[
\widetilde P|_{D(f)}
\cong
\mathcal O_{D(f)}^{\oplus r},
\]
and \(\widetilde P\) is finite locally free.

Conversely, let \(\mathcal E\) be finite locally free on \(X\).  Since it is quasi-coherent,
\[
\mathcal E\cong\widetilde P,
\qquad
P=\Gamma(X,\mathcal E).
\]
Finite local freeness gives a cover by distinguished opens
\[
D(f_i)
\]
such that
\[
P_{f_i}\cong A_{f_i}^{\oplus r_i}.
\]
Because \(X\) is quasi-compact, choose finitely many such opens.  Hence \(P\) is locally finite free at every prime
and finitely presented.  A finitely presented module whose localizations at all primes are free is projective.
Thus \(P\) is finitely generated projective.

Morphisms correspond as well:
\[
\operatorname{Hom}_{\mathcal O_X}(\widetilde P,\widetilde Q)
\cong
\operatorname{Hom}_A(P,Q).
\]
Therefore
\[
\begin{gathered}
\{\text{finitely generated projective }A\text{-modules}\}\\
\simeq\\
\{\text{finite locally free sheaves on }\Spec A\}.
\end{gathered}
\]

Now define the rank function
\[
\rho_P:\Spec A\to\mathbb Z_{\ge0},
\qquad
\rho_P(\mathfrak p)
=
\dim_{\kappa(\mathfrak p)}
(P\otimes_A\kappa(\mathfrak p)).
\]
On any distinguished open where \(P\) is free of rank \(r\), the function is constantly \(r\). Hence \(\rho_P\)
is locally constant.  Therefore it is constant on every connected component of \(\Spec A\).

This explains the geometric language: a finite locally free sheaf has locally constant rank, and on a connected
affine scheme a finitely generated projective module has a single constant rank.
\end{solution}

% VI33-DETAILED-CHALLENGE-03
\subsection*{Challenge VI.33.C3 solution}
\begin{solution}
Take
\[
A=\mathbb Z
\qquad\text{and}\qquad
M=\mathbb Q/\mathbb Z.
\]
Every element of \(M\) is torsion, but there is no single nonzero integer killing all of \(M\). Therefore
\[
\operatorname{Ann}_{\mathbb Z}(M)=0.
\]
Hence
\[
V(\operatorname{Ann}M)=V(0)=\Spec\mathbb Z.
\]

We now compute the support.  At the generic prime \((0)\), localization inverts every nonzero integer.  Given
\[
m\in\mathbb Q/\mathbb Z,
\]
some nonzero integer \(n\) kills \(m\).  Since \(n\) becomes a unit in \(\mathbb Z_{(0)}=\mathbb Q\), every
localized element vanishes. Thus
\[
M_{(0)}=0.
\]

At a closed prime \((p)\), the element
\[
\frac1p+\mathbb Z
\]
is nonzero after localization at \((p)\).  Indeed an integer \(s\notin(p)\) is prime to \(p\), and multiplication
by \(s\) cannot kill this element of order \(p\). Hence
\[
M_{(p)}\ne0
\]
for every prime number \(p\).

Therefore
\[
\boxed{
\operatorname{Supp}(\widetilde M)
=
\{(p):p\text{ prime}\},
}
\]
while
\[
\boxed{
V(\operatorname{Ann}M)
=
\Spec\mathbb Z
=
\{(0)\}\cup\{(p):p\text{ prime}\}.
}
\]
So the support is not even closed.

Where did finite generation enter the usual proof?  If a finite generating set
\[
m_1,\ldots,m_r
\]
all vanish in \(M_{\mathfrak p}\), one chooses
\[
s_i\notin\mathfrak p
\quad\text{with}\quad
s_i m_i=0
\]
and multiplies the finitely many \(s_i\) to obtain one element
\[
s=s_1\cdots s_r\notin\mathfrak p
\]
annihilating all of \(M\).  For \(\mathbb Q/\mathbb Z\), every individual element has an annihilator outside
\((0)\), but there is no finite generating set from which one common annihilator can be formed.
\end{solution}

% VI33-DETAILED-CHALLENGE-04
\subsection*{Challenge VI.33.C4 solution}
\begin{solution}
Let
\[
A\to B
\]
and let \(I\subseteq A\).  The closed subscheme
\[
Z=\Spec(A/I)
\hookrightarrow
X=\Spec A
\]
has base change
\[
Z_B
=
Z\times_X\Spec B.
\]
The affine fiber-product formula gives
\[
Z_B
=
\Spec\bigl((A/I)\otimes_A B\bigr).
\]
But
\[
(A/I)\otimes_A B
\cong
B/IB,
\]
so
\[
\boxed{
Z_B\cong\Spec(B/IB).
}
\]

Now recover this from the ideal sheaf.  On \(X\) we have
\[
0\to\widetilde I\to\mathcal O_X\to i_*\mathcal O_Z\to0.
\]
Pullback is tensorial and therefore right exact, but it is \emph{not} left exact in general.  Thus one must not
blindly write a new short exact sequence beginning with \(0\).

Instead,
\[
f^*\widetilde I
\cong
\widetilde{B\otimes_A I}.
\]
The pulled-back morphism
\[
f^*\widetilde I\longrightarrow\mathcal O_{\Spec B}
\]
corresponds on modules to the multiplication map
\[
B\otimes_A I\longrightarrow B,
\qquad
b\otimes i\longmapsto b\,\varphi(i).
\]
Its image is exactly the extended ideal
\[
IB\subseteq B.
\]
Therefore the image sheaf is
\[
\widetilde{IB}\subseteq\mathcal O_{\Spec B}.
\]
The cokernel is
\[
\widetilde{B/IB}.
\]
Hence the closed subscheme cut out after base change is
\[
\boxed{\Spec(B/IB)}.
\]

This distinction is important: \(B\otimes_A I\to B\) need not be injective when \(B\) is not flat over \(A\).
The ideal of the pullback closed subscheme is the \emph{image} \(IB\), not an unjustified identification of
\(B\otimes_A I\) with a submodule of \(B\).
\end{solution}

% VI33-DETAILED-CHALLENGE-05
\subsection*{Challenge VI.33.C5 solution}
\begin{solution}
Let
\[
S=\bigoplus_{n\ge0}S_n
\]
be a graded ring and let
\[
M=\bigoplus_{n\in\mathbb Z}M_n
\]
be a graded \(S\)-module.  The affine construction suggests that the sheaf should be determined on a basis of
standard opens by localization.

On \(\Proj S\), the standard opens are
\[
D_+(f)
=
\{\mathfrak p\in\Proj S:f\notin\mathfrak p\}
\]
for homogeneous \(f\in S_+\).  The affine coordinate ring of \(D_+(f)\) is
\[
S_{(f)}
:=
(S_f)_0,
\]
the degree-zero part of the homogeneous localization.

Accordingly define the local module
\[
M_{(f)}
:=
(M_f)_0.
\]
It is naturally an \(S_{(f)}\)-module.  The proposed sheaf recipe is therefore
\[
\boxed{
D_+(f)\longmapsto M_{(f)}.
}
\]

We must check overlap compatibility.  If \(f\) and \(g\) are homogeneous, then
\[
D_+(f)\cap D_+(g)=D_+(fg).
\]
Restricting the \(S_{(f)}\)-module \(M_{(f)}\) to the overlap means localizing further by an appropriate
degree-zero element obtained from a power of \(g\).  The result canonically identifies with
\[
M_{(fg)}.
\]
The same module is obtained by first starting from \(M_{(g)}\) and restricting.  Associativity of homogeneous
localization gives the cocycle compatibility on triple overlaps.

Thus the \(M_{(f)}\) glue to an \(\mathcal O_{\Proj S}\)-module, customarily denoted again by
\[
\widetilde M.
\]

The essential difference from the affine case is therefore not the sheaf logic but the grading:
\[
\boxed{
\text{affine: }M_f
\qquad\longrightarrow\qquad
\text{projective: }(M_f)_0.
}
\]
This predicts the later construction of twisting sheaves as well: shifting the grading of \(M\) changes the
degree-zero pieces after localization, even though the underlying ungraded module may look similar.
\end{solution}
